\documentclass[12pt,a4paper]{article}

\usepackage[margin=1in]{geometry}
\usepackage{amsmath, amssymb, amsthm}
\usepackage{enumitem}
\usepackage{booktabs}

\setlength{\parindent}{0pt}
\setlength{\parskip}{6pt}

\newtheoremstyle{notestyle}
  {8pt}
  {8pt}
  {\normalfont}
  {}
  {\bfseries}
  {.}
  {0.5em}
  {}

\theoremstyle{notestyle}
\newtheorem{definition}{Definition}
\newtheorem{example}{Example}

\newenvironment{answer}
{\par\textbf{Answer.}\ }
{\hfill$\square$\par}

\title{\textbf{Supplementary Practice: Stationary Distribution Calculations}}
\author{MA4034: Time Series and Stochastic Processes}
\date{}

\begin{document}
\maketitle

\section{Main Idea}

For a Markov chain with transition matrix $P$, a stationary distribution is a probability vector
\[
\pi=(\pi_0,\pi_1,\ldots)
\]
such that
\[
\pi P=\pi,
\]
where
\[
\pi_i\geq 0
\qquad\text{and}\qquad
\sum_i \pi_i=1.
\]

The equation
\[
\pi P=\pi
\]
means that if the chain starts with distribution $\pi$, then after one step the distribution is still $\pi$.

\textbf{Interpretation.} In an ergodic Markov chain, the stationary distribution gives the long-run proportion of time spent in each state.

\section{How to Prove Whether a Stationary Distribution Exists}

When a question asks whether a Markov chain has a stationary distribution, use the following steps.

\subsection*{Case 1: Finite State Space}

If the Markov chain has a finite number of states, a stationary distribution always exists.

However, it may or may not be unique.

\begin{itemize}[leftmargin=1.2cm]
    \item If the finite chain is irreducible, then the stationary distribution exists and is unique.
    \item If the finite chain is not irreducible, then a stationary distribution still exists, but it may not be unique.
\end{itemize}

\textbf{Steps for a finite chain:}

\begin{enumerate}[leftmargin=1.2cm]
    \item Write
    \[
    \pi=(\pi_0,\pi_1,\ldots,\pi_m).
    \]
    \item Solve
    \[
    \pi P=\pi.
    \]
    \item Add the probability condition
    \[
    \pi_0+\pi_1+\cdots+\pi_m=1.
    \]
    \item Check that all probabilities are non-negative:
    \[
    \pi_i\geq 0.
    \]
    \item If one valid probability vector is obtained, then a stationary distribution exists.
\end{enumerate}

\subsection*{Case 2: Infinite State Space}

For an infinite state space, a stationary distribution may or may not exist. We must be more careful.

\textbf{Steps for an infinite chain:}

\begin{enumerate}[leftmargin=1.2cm]
    \item Write
    \[
    \pi=(\pi_0,\pi_1,\pi_2,\ldots).
    \]
    \item Solve the equation
    \[
    \pi P=\pi
    \]
    to express $\pi_1,\pi_2,\ldots$ in terms of $\pi_0$.
    \item Use the probability condition
    \[
    \sum_{i=0}^{\infty}\pi_i=1.
    \]
    \item Check whether the infinite series converges.
    \item If the series converges and all $\pi_i\geq 0$, then a stationary distribution exists.
    \item If the series diverges, then no stationary distribution exists.
\end{enumerate}

\subsection*{Case 3: Using Irreducibility and Recurrence}

For an irreducible Markov chain:

\[
\boxed{
\text{A stationary distribution exists if and only if the chain is positive recurrent.}
}
\]

So for an irreducible chain:

\begin{itemize}[leftmargin=1.2cm]
    \item If the chain is positive recurrent, then a unique stationary distribution exists.
    \item If the chain is null recurrent or transient, then no stationary distribution exists.
\end{itemize}

\subsection*{Case 4: If a Candidate Distribution is Given}

If the question gives a vector and asks whether it is stationary:

\begin{enumerate}[leftmargin=1.2cm]
    \item Check whether all entries are non-negative.
    \item Check whether the entries sum to 1.
    \item Compute
    \[
    \pi P.
    \]
    \item If
    \[
    \pi P=\pi,
    \]
    then $\pi$ is stationary.
    \item If
    \[
    \pi P\neq \pi,
    \]
    then $\pi$ is not stationary.
\end{enumerate}

\subsection*{Common Exam Mistake}

Do not stop after solving
\[
\pi P=\pi.
\]
You must also check
\[
\sum_i\pi_i=1
\]
and
\[
\pi_i\geq 0.
\]

A vector can satisfy the algebraic equation but still fail to be a probability distribution.

\subsection*{Quick Decision Table}

\[
\begin{array}{c|c}
\text{Situation} & \text{Conclusion}\\
\hline
\text{Finite and irreducible} & \text{unique stationary distribution exists}\\
\text{Finite but reducible} & \text{stationary distribution exists, may not be unique}\\
\text{Infinite and irreducible positive recurrent} & \text{unique stationary distribution exists}\\
\text{Infinite and irreducible null recurrent} & \text{no stationary distribution}\\
\text{Infinite and irreducible transient} & \text{no stationary distribution}\\
\text{Candidate }\pi\text{ given} & \text{check }\pi_i\geq0,\ \sum_i\pi_i=1,\ \pi P=\pi
\end{array}
\]

\section{Worked Calculation 1: Two-State Chain}

\begin{example}
Consider a Markov chain with transition matrix
\[
P=
\begin{pmatrix}
0.7 & 0.3\\
0.4 & 0.6
\end{pmatrix}.
\]
Find the stationary distribution.
\end{example}

\begin{answer}
Let
\[
\pi=(\pi_0,\pi_1).
\]
We solve
\[
\pi P=\pi
\]
and
\[
\pi_0+\pi_1=1.
\]

Now,
\[
(\pi_0,\pi_1)
\begin{pmatrix}
0.7 & 0.3\\
0.4 & 0.6
\end{pmatrix}
=
(\pi_0,\pi_1).
\]

Equating the first component:
\[
0.7\pi_0+0.4\pi_1=\pi_0.
\]
Rearrange:
\[
0.4\pi_1=0.3\pi_0.
\]
Therefore,
\[
\pi_1=\frac{0.3}{0.4}\pi_0=\frac34\pi_0.
\]

Using the probability condition,
\[
\pi_0+\pi_1=1,
\]
we get
\[
\pi_0+\frac34\pi_0=1.
\]
Thus,
\[
\frac74\pi_0=1,
\]
so
\[
\pi_0=\frac47.
\]
Then
\[
\pi_1=\frac34\cdot\frac47=\frac37.
\]

Therefore the stationary distribution is
\[
\boxed{\pi=\left(\frac47,\frac37\right).}
\]
\end{answer}

\section{Worked Calculation 2: Weather Chain}

\begin{example}
Suppose the weather is classified as sunny $(S)$ or rainy $(R)$. A sunny day is followed by a sunny day with probability $0.8$, and a rainy day is followed by a sunny day with probability $0.5$.

Using the state order $(S,R)$, find the stationary distribution.
\end{example}

\begin{answer}
First construct the transition matrix.

If today is sunny:
\[
\Prob(S\to S)=0.8,
\qquad
\Prob(S\to R)=0.2.
\]

If today is rainy:
\[
\Prob(R\to S)=0.5,
\qquad
\Prob(R\to R)=0.5.
\]

Therefore,
\[
P=
\begin{pmatrix}
0.8 & 0.2\\
0.5 & 0.5
\end{pmatrix}.
\]

Let
\[
\pi=(\pi_S,\pi_R).
\]
We solve
\[
\pi P=\pi
\]
and
\[
\pi_S+\pi_R=1.
\]

The first component gives
\[
0.8\pi_S+0.5\pi_R=\pi_S.
\]
Rearrange:
\[
0.5\pi_R=0.2\pi_S.
\]
So
\[
\pi_R=\frac{0.2}{0.5}\pi_S=\frac25\pi_S.
\]

Using
\[
\pi_S+\pi_R=1,
\]
we get
\[
\pi_S+\frac25\pi_S=1.
\]
Thus,
\[
\frac75\pi_S=1,
\]
so
\[
\pi_S=\frac57.
\]
Then
\[
\pi_R=\frac27.
\]

Therefore,
\[
\boxed{\pi=\left(\frac57,\frac27\right).}
\]

Interpretation: in the long run, about $\frac57$ of days are sunny and $\frac27$ of days are rainy.
\end{answer}

\section{Worked Calculation 3: Three-State Chain}

\begin{example}
Consider a Markov chain with state space
\[
S=\{0,1,2\}
\]
and transition matrix
\[
P=
\begin{pmatrix}
0.5 & 0.5 & 0\\
0.2 & 0.5 & 0.3\\
0 & 0.4 & 0.6
\end{pmatrix}.
\]
Find the stationary distribution.
\end{example}

\begin{answer}
Let
\[
\pi=(\pi_0,\pi_1,\pi_2).
\]
We solve
\[
\pi P=\pi
\]
and
\[
\pi_0+\pi_1+\pi_2=1.
\]

Writing out $\pi P=\pi$:
\[
0.5\pi_0+0.2\pi_1=\pi_0,
\]
\[
0.5\pi_0+0.5\pi_1+0.4\pi_2=\pi_1,
\]
\[
0.3\pi_1+0.6\pi_2=\pi_2.
\]

Use the first equation:
\[
0.5\pi_0+0.2\pi_1=\pi_0.
\]
Then
\[
0.2\pi_1=0.5\pi_0,
\]
so
\[
\pi_1=\frac{0.5}{0.2}\pi_0=\frac52\pi_0.
\]

Use the third equation:
\[
0.3\pi_1+0.6\pi_2=\pi_2.
\]
Then
\[
0.3\pi_1=0.4\pi_2,
\]
so
\[
\pi_2=\frac{0.3}{0.4}\pi_1=\frac34\pi_1.
\]

Substitute $\pi_1=\frac52\pi_0$:
\[
\pi_2=\frac34\cdot\frac52\pi_0=\frac{15}{8}\pi_0.
\]

Now use the sum condition:
\[
\pi_0+\pi_1+\pi_2=1.
\]
Substitute everything in terms of $\pi_0$:
\[
\pi_0+\frac52\pi_0+\frac{15}{8}\pi_0=1.
\]
Use denominator 8:
\[
\frac88\pi_0+\frac{20}{8}\pi_0+\frac{15}{8}\pi_0=1.
\]
Therefore,
\[
\frac{43}{8}\pi_0=1.
\]
So
\[
\pi_0=\frac8{43}.
\]

Then
\[
\pi_1=\frac52\cdot\frac8{43}=\frac{20}{43},
\]
and
\[
\pi_2=\frac{15}{8}\cdot\frac8{43}=\frac{15}{43}.
\]

Therefore the stationary distribution is
\[
\boxed{
\pi=\left(\frac8{43},\frac{20}{43},\frac{15}{43}\right).
}
\]
\end{answer}

\section{Practice Questions}

\begin{enumerate}[label=\textbf{Q\arabic*.}, leftmargin=1.2cm]
    \item Find the stationary distribution of
    \[
    P=
    \begin{pmatrix}
    0.6 & 0.4\\
    0.2 & 0.8
    \end{pmatrix}.
    \]

    \item A machine has two states: working $(W)$ and broken $(B)$. If it is working today, it remains working tomorrow with probability $0.9$. If it is broken today, it is repaired and working tomorrow with probability $0.6$. Find the long-run proportion of time the machine is working.

    \item Find the stationary distribution of
    \[
    P=
    \begin{pmatrix}
    0.4 & 0.6 & 0\\
    0.3 & 0.4 & 0.3\\
    0 & 0.2 & 0.8
    \end{pmatrix}.
    \]

    \item For the chain
    \[
    P=
    \begin{pmatrix}
    0.7 & 0.3\\
    0.1 & 0.9
    \end{pmatrix},
    \]
    find the stationary distribution and interpret it in words.
\end{enumerate}

\section{Answers}

\subsection*{Answer 1}

Let
\[
\pi=(\pi_0,\pi_1).
\]
Solve
\[
\pi P=\pi,
\qquad
\pi_0+\pi_1=1.
\]

From the first component:
\[
0.6\pi_0+0.2\pi_1=\pi_0.
\]
Thus,
\[
0.2\pi_1=0.4\pi_0,
\]
so
\[
\pi_1=2\pi_0.
\]

Using
\[
\pi_0+\pi_1=1,
\]
we get
\[
\pi_0+2\pi_0=1.
\]
Therefore,
\[
\pi_0=\frac13
\]
and
\[
\pi_1=\frac23.
\]

Hence,
\[
\boxed{\pi=\left(\frac13,\frac23\right).}
\]

\subsection*{Answer 2}

Use the state order $(W,B)$.

If the machine is working:
\[
\Prob(W\to W)=0.9,
\qquad
\Prob(W\to B)=0.1.
\]

If the machine is broken:
\[
\Prob(B\to W)=0.6,
\qquad
\Prob(B\to B)=0.4.
\]

Thus,
\[
P=
\begin{pmatrix}
0.9 & 0.1\\
0.6 & 0.4
\end{pmatrix}.
\]

Let
\[
\pi=(\pi_W,\pi_B).
\]
From $\pi P=\pi$,
\[
0.9\pi_W+0.6\pi_B=\pi_W.
\]
Therefore,
\[
0.6\pi_B=0.1\pi_W.
\]
So
\[
\pi_B=\frac{1}{6}\pi_W.
\]

Using
\[
\pi_W+\pi_B=1,
\]
we get
\[
\pi_W+\frac16\pi_W=1.
\]
Thus,
\[
\frac76\pi_W=1,
\]
so
\[
\pi_W=\frac67.
\]

Therefore the long-run proportion of time the machine is working is
\[
\boxed{\frac67\approx 0.8571.}
\]

\subsection*{Answer 3}

Let
\[
\pi=(\pi_0,\pi_1,\pi_2).
\]

From $\pi P=\pi$:
\[
0.4\pi_0+0.3\pi_1=\pi_0,
\]
\[
0.6\pi_0+0.4\pi_1+0.2\pi_2=\pi_1,
\]
\[
0.3\pi_1+0.8\pi_2=\pi_2.
\]

From the first equation:
\[
0.3\pi_1=0.6\pi_0,
\]
so
\[
\pi_1=2\pi_0.
\]

From the third equation:
\[
0.3\pi_1=0.2\pi_2,
\]
so
\[
\pi_2=\frac{0.3}{0.2}\pi_1=\frac32\pi_1.
\]
Since $\pi_1=2\pi_0$,
\[
\pi_2=3\pi_0.
\]

Use
\[
\pi_0+\pi_1+\pi_2=1.
\]
Then
\[
\pi_0+2\pi_0+3\pi_0=1.
\]
So
\[
6\pi_0=1,
\]
and
\[
\pi_0=\frac16.
\]
Therefore,
\[
\pi_1=\frac13,
\qquad
\pi_2=\frac12.
\]

Hence,
\[
\boxed{\pi=\left(\frac16,\frac13,\frac12\right).}
\]

\subsection*{Answer 4}

Let
\[
\pi=(\pi_0,\pi_1).
\]

From the first component of $\pi P=\pi$:
\[
0.7\pi_0+0.1\pi_1=\pi_0.
\]
Thus,
\[
0.1\pi_1=0.3\pi_0,
\]
so
\[
\pi_1=3\pi_0.
\]

Using
\[
\pi_0+\pi_1=1,
\]
we get
\[
\pi_0+3\pi_0=1.
\]
Therefore,
\[
\pi_0=\frac14,
\qquad
\pi_1=\frac34.
\]

Thus,
\[
\boxed{\pi=\left(\frac14,\frac34\right).}
\]

Interpretation: in the long run, the chain spends one quarter of the time in state 0 and three quarters of the time in state 1.

\section{Checking Whether a Given Distribution is Stationary}

Sometimes the question does not ask us to find the stationary distribution. Instead, it gives a probability vector and asks whether it is stationary.

To check this, use two steps:
\begin{enumerate}[leftmargin=1.2cm]
    \item Check whether the vector is a probability distribution:
    \[
    \pi_i\geq 0
    \quad\text{and}\quad
    \sum_i\pi_i=1.
    \]
    \item Check whether
    \[
    \pi P=\pi.
    \]
\end{enumerate}

If both conditions are true, then $\pi$ is stationary. If either condition fails, then $\pi$ is not stationary.

\begin{example}
Consider
\[
P=
\begin{pmatrix}
0.6 & 0.4\\
0.2 & 0.8
\end{pmatrix}.
\]
Is
\[
\pi=\left(\frac13,\frac23\right)
\]
a stationary distribution? Give a reason.
\end{example}

\begin{answer}
First check that $\pi$ is a probability distribution:
\[
\frac13+\frac23=1,
\]
and both entries are non-negative.

Now calculate $\pi P$:
\[
\pi P
=
\left(\frac13,\frac23\right)
\begin{pmatrix}
0.6 & 0.4\\
0.2 & 0.8
\end{pmatrix}.
\]

First component:
\[
\frac13(0.6)+\frac23(0.2)
=
\frac{0.6}{3}+\frac{0.4}{3}
=
\frac{1}{3}.
\]

Second component:
\[
\frac13(0.4)+\frac23(0.8)
=
\frac{0.4}{3}+\frac{1.6}{3}
=
\frac{2}{3}.
\]

Therefore,
\[
\pi P=\left(\frac13,\frac23\right)=\pi.
\]

Hence,
\[
\boxed{\pi=\left(\frac13,\frac23\right)\text{ is stationary.}}
\]
\end{answer}

\section{Practice Questions: Is the Given Vector Stationary?}

\begin{enumerate}[label=\textbf{Q\arabic*.}, leftmargin=1.2cm]
    \item Consider
    \[
    P=
    \begin{pmatrix}
    0.7 & 0.3\\
    0.4 & 0.6
    \end{pmatrix}.
    \]
    Is
    \[
    \pi=\left(\frac47,\frac37\right)
    \]
    a stationary distribution? Give a reason.

    \item Consider
    \[
    P=
    \begin{pmatrix}
    0.7 & 0.3\\
    0.4 & 0.6
    \end{pmatrix}.
    \]
    Is
    \[
    \pi=\left(\frac12,\frac12\right)
    \]
    a stationary distribution? Give a reason.

    \item Consider
    \[
    P=
    \begin{pmatrix}
    1 & 0 & 0\\
    \frac14 & \frac12 & \frac14\\
    0 & 0 & 1
    \end{pmatrix}.
    \]
    Is
    \[
    \pi=\left(\frac12,0,\frac12\right)
    \]
    a stationary distribution? Give a reason.

    \item For the same matrix
    \[
    P=
    \begin{pmatrix}
    1 & 0 & 0\\
    \frac14 & \frac12 & \frac14\\
    0 & 0 & 1
    \end{pmatrix},
    \]
    is
    \[
    \pi=\left(0,\frac12,\frac12\right)
    \]
    a stationary distribution? Give a reason.

    \item Consider
    \[
    P=
    \begin{pmatrix}
    0.5 & 0.5\\
    0.1 & 0.9
    \end{pmatrix}.
    \]
    Is
    \[
    v=(2,-1)
    \]
    a stationary distribution? Give a reason.
\end{enumerate}

\section{Answers: Is the Given Vector Stationary?}

\subsection*{Answer 1}

First,
\[
\frac47+\frac37=1,
\]
so $\pi$ is a probability distribution.

Now compute:
\[
\pi P
=
\left(\frac47,\frac37\right)
\begin{pmatrix}
0.7 & 0.3\\
0.4 & 0.6
\end{pmatrix}.
\]

First component:
\[
\frac47(0.7)+\frac37(0.4)
=
\frac47\cdot\frac7{10}+\frac37\cdot\frac4{10}
=
\frac25+\frac6{35}
=
\frac{14}{35}+\frac6{35}
=
\frac{20}{35}
=
\frac47.
\]

Second component:
\[
\frac47(0.3)+\frac37(0.6)
=
\frac47\cdot\frac3{10}+\frac37\cdot\frac6{10}
=
\frac6{35}+\frac9{35}
=
\frac{15}{35}
=
\frac37.
\]

Thus,
\[
\pi P=\left(\frac47,\frac37\right)=\pi.
\]

Therefore,
\[
\boxed{\pi\text{ is stationary.}}
\]

\subsection*{Answer 2}

First,
\[
\frac12+\frac12=1,
\]
so $\pi$ is a probability distribution.

Now compute:
\[
\pi P
=
\left(\frac12,\frac12\right)
\begin{pmatrix}
0.7 & 0.3\\
0.4 & 0.6
\end{pmatrix}.
\]

First component:
\[
\frac12(0.7)+\frac12(0.4)
=0.35+0.20
=0.55.
\]

Second component:
\[
\frac12(0.3)+\frac12(0.6)
=0.15+0.30
=0.45.
\]

Therefore,
\[
\pi P=(0.55,0.45).
\]

But
\[
(0.55,0.45)\neq \left(\frac12,\frac12\right).
\]

Therefore,
\[
\boxed{\pi\text{ is not stationary.}}
\]

\subsection*{Answer 3}

First,
\[
\frac12+0+\frac12=1,
\]
so $\pi$ is a probability distribution.

Now compute:
\[
\pi P
=
\left(\frac12,0,\frac12\right)
\begin{pmatrix}
1 & 0 & 0\\
\frac14 & \frac12 & \frac14\\
0 & 0 & 1
\end{pmatrix}.
\]

First component:
\[
\frac12(1)+0\left(\frac14\right)+\frac12(0)=\frac12.
\]

Second component:
\[
\frac12(0)+0\left(\frac12\right)+\frac12(0)=0.
\]

Third component:
\[
\frac12(0)+0\left(\frac14\right)+\frac12(1)=\frac12.
\]

Thus,
\[
\pi P=\left(\frac12,0,\frac12\right)=\pi.
\]

Therefore,
\[
\boxed{\pi\text{ is stationary.}}
\]

\subsection*{Answer 4}

First,
\[
0+\frac12+\frac12=1,
\]
so $\pi$ is a probability distribution.

Now compute:
\[
\pi P
=
\left(0,\frac12,\frac12\right)
\begin{pmatrix}
1 & 0 & 0\\
\frac14 & \frac12 & \frac14\\
0 & 0 & 1
\end{pmatrix}.
\]

First component:
\[
0(1)+\frac12\left(\frac14\right)+\frac12(0)=\frac18.
\]

Second component:
\[
0(0)+\frac12\left(\frac12\right)+\frac12(0)=\frac14.
\]

Third component:
\[
0(0)+\frac12\left(\frac14\right)+\frac12(1)=\frac18+\frac12=\frac58.
\]

Therefore,
\[
\pi P=\left(\frac18,\frac14,\frac58\right).
\]

But
\[
\left(\frac18,\frac14,\frac58\right)
\neq
\left(0,\frac12,\frac12\right).
\]

Therefore,
\[
\boxed{\pi\text{ is not stationary.}}
\]

\subsection*{Answer 5}

A stationary distribution must be a probability distribution. That means all entries must be non-negative and must sum to 1.

For
\[
v=(2,-1),
\]
the second entry is negative:
\[
-1<0.
\]

Although
\[
2+(-1)=1,
\]
the vector has a negative entry, so it is not a probability distribution.

Therefore it cannot be a stationary distribution.

\[
\boxed{v\text{ is not stationary.}}
\]

\end{document}
